\documentclass[11pt]{article}
\usepackage[utf8]{inputenc}
\usepackage{amsmath,amssymb,amsthm}
\usepackage[margin=1in]{geometry}
\usepackage{hyperref}
\usepackage{xcolor}
\usepackage{enumitem}
\usepackage{newunicodechar}
\newunicodechar{ℝ}{\ensuremath{\mathbb{R}}}
\newunicodechar{ℕ}{\ensuremath{\mathbb{N}}}
\newunicodechar{⟨}{\ensuremath{\langle}}
\newunicodechar{⟩}{\ensuremath{\rangle}}
\newunicodechar{α}{\ensuremath{\alpha}}
\newunicodechar{π}{\ensuremath{\pi}}
\newunicodechar{Γ}{\ensuremath{\Gamma}}
\newunicodechar{≤}{\ensuremath{\le}}
\newunicodechar{·}{\ensuremath{\cdot}}
\newunicodechar{→}{\ensuremath{\to}}

\newcommand{\userbox}[1]{\par\medskip\begin{quote}\small\textbf{\textcolor{blue!60}{DM:}} #1\end{quote}\medskip}
\newcommand{\claudebox}[1]{\par\medskip\begin{quote}\small\textbf{\textcolor{orange!60!black}{Claude:}} #1\end{quote}\medskip}
\newcommand{\gptbox}[1]{\par\medskip\begin{quote}\small\textbf{\textcolor{green!50!black}{GPT-5.5 Pro:}} #1\end{quote}\medskip}

\title{Tide retrospective:\\\emph{generic-$k$ J-function asymptotic
for monomial potentials}}
\author{Claude Opus 4.7 (1M ctx), with Daniel Murfet}
\date{2026-05-07}

\begin{document}
\maketitle

\begin{abstract}
For the single-monomial potential $K(w) = w^{2k}/(2k)!$ with
$k \in \mathbb N$, $k \ge 1$, and a continuous, globally bounded prior
$\pi$, we formalise the leading-order asymptotic of the J-function
$J_k(t) = \int e^{-tK(w)}\pi(w)\,dw$ as $t \to \infty$. The headline
constant is
$\pi(0)\cdot \tfrac{1}{k}\cdot ((2k)!)^{1/(2k)}\cdot \Gamma(1/(2k))$.
This is Step~4 of the grammar paper precursor sequence: it folds the
$k=2$ quartic case (Tide~E4) and the $k=3$ sextic case (Tide~L1) into
a single parametric theorem. New file
\texttt{Laplace/OneD/JFunctionMonomial.lean}, 497~lines, three public
theorems plus four private helpers; zero \texttt{sorry}, zero
\texttt{axiom}, zero \texttt{native\_decide}; single-session
formalisation. The Step~2 deliberation with GPT-5.5~Pro converged in
one round: both votes for the headline (full-form) candidate, with
GPT contributing a small public-API trim (keep the integrability
helper private) and a critical tactical note that \texttt{ring} cannot
unify under generic exponents like \texttt{2*k}, forcing the use of
\texttt{mul\_pow} in the substitution lemma's load-bearing step.
\end{abstract}

\section{Setting}

The grammar paper precursor sequence is a chain of formalised
asymptotic-method results that ground the grammar paper's Sections
on resolution of singularities. Three earlier tides have laid the
seabed:

\begin{itemize}[leftmargin=1em,itemsep=0.2em]
\item Tide~7 (the original \emph{grammar precursor} tide, 2~May)
  established the substitute-then-DCT skeleton.
\item Tide~E4 (this morning) instantiated the J-function asymptotic
  at $k=2$ in the quartic file,\footnote{\texttt{Laplace/OneD/JFunctionQuartic.lean}.}
  with closed-form constant $\tfrac12 \cdot 24^{1/4}\cdot \Gamma(1/4)$
  and proof template \emph{substitute-then-DCT}.
\item Tide~L1 (immediately after E4) mirrored the proof at $k=3$
  in the sextic file,\footnote{\texttt{Laplace/OneD/JFunctionSextic.lean}.}
  with constant $\tfrac13\cdot 720^{1/6}\cdot \Gamma(1/6)$.
\end{itemize}

\noindent The L1 retrospective explicitly named the generic-$k$
abstraction as \emph{``now ripe''}: with both $k=2$ and $k=3$ in the
same shape, the abstraction over $K(w) = w^{2k}/(2k)!$ becomes the
natural next step. The constant generalises cleanly to
$\tfrac{1}{k}\cdot ((2k)!)^{1/(2k)}\cdot \Gamma(1/(2k))$, and the
proof template generalises modulo cast bookkeeping.

The Step~2 deliberation between Claude and GPT-5.5~Pro picked the
headline candidate (full asymptotic plus closed-form partition, three
public theorems) over a thinner alternative that would have postponed
the asymptotic corollary as a follow-up. Both models voted the same
target in a single round; GPT additionally trimmed the public API
(\texttt{kth\_integrable\_pow} kept private, since only $n=0$ is
needed inside this file) and reshaped the proof tactics to use
\texttt{mul\_pow} rather than the \texttt{ring} idiom from L1
(\texttt{ring} cannot unify under a parametric \texttt{2*k}).

The markdown tide log\footnote{In the SRI repo at
\texttt{projects/primer/tide-log/}, file
\texttt{2026-05-07-tide-generic-k-j-function.md}.}
records the full deliberation; the GPT consult is preserved at the
sibling \texttt{gpt\_responses\_n1/n1\_v1.md}.

\section{The thing formalised}

\begin{quote}\itshape
For $k \in \mathbb N$ with $k \ge 1$ and a continuous, globally
bounded prior $\pi : \mathbb R \to \mathbb R$ (with global bound $M$),
$$
  \lim_{t \to \infty} t^{1/(2k)} \cdot
    \int_{\mathbb R} e^{-t\,w^{2k}/(2k)!}\,\pi(w)\,dw
  \;=\; \pi(0) \cdot \tfrac{1}{k} \cdot ((2k)!)^{1/(2k)}
        \cdot \Gamma\!\left(\tfrac{1}{2k}\right).
$$
\end{quote}

\noindent The Lean signature\footnote{\texttt{Laplace/OneD/JFunctionMonomial.lean},
\texttt{kth\_jfunction\_asymptotic}.}:

\begin{verbatim}
theorem kth_jfunction_asymptotic
    {k : ℕ} (hk : 1 ≤ k)
    {prior : ℝ → ℝ} (hprior_cont : Continuous prior)
    {M : ℝ} (hprior_bd : ∀ x, |prior x| ≤ M) :
    Tendsto (fun t : ℝ =>
        t ^ ((1 : ℝ) / ((2 * k : ℕ) : ℝ)) *
        ∫ w : ℝ, Real.exp (-(t * w ^ (2 * k) /
                              (Nat.factorial (2 * k) : ℝ)))
                  * prior w)
      atTop
      (𝓝 (prior 0 * (1 / (k : ℝ))
            * ((Nat.factorial (2 * k) : ℝ))
                ^ ((1 : ℝ) / ((2 * k : ℕ) : ℝ))
            * Real.Gamma ((1 : ℝ) / ((2 * k : ℕ) : ℝ))))
```

\noindent The two existing files (\texttt{JFunctionQuartic.lean} for
$k=2$, \texttt{JFunctionSextic.lean} for $k=3$) are not yet refactored
as 5-line specialisations of this theorem, since both are on unmerged
branches at the base time; that is a follow-up tide once all three
land on \texttt{main}.

The closed-form partition is also exported:

\begin{verbatim}
theorem kth_partition {k : ℕ} (hk : 1 ≤ k) {t : ℝ} (ht : 0 < t) :
    ∫ x : ℝ, Real.exp (-(t * x ^ (2 * k) /
                            (Nat.factorial (2 * k) : ℝ)))
      = (1 / (k : ℝ))
          * ((Nat.factorial (2 * k) : ℝ) / t)
              ^ ((1 : ℝ) / ((2 * k : ℕ) : ℝ))
          * Real.Gamma ((1 : ℝ) / ((2 * k : ℕ) : ℝ))
\end{verbatim}

\noindent and the centred DCT core
\texttt{kth\_jfunction\_centered\_tendsto\_zero}, which is the
load-bearing intermediate.

The hypothesis $k \ge 1$ is mathematically essential: at $k=0$ the
``potential'' $w^{0}/0! = 1$ is constant, making
$e^{-t}\,\pi(w)$ have no concentration behaviour and the asymptotic
trivial in the wrong way. Continuous-and-bounded on the prior is the
slightly stronger form taken in E4/L1 over the deliberation's
\texttt{AEMeasurable + ContinuousAt 0}; weakening across all three
files is a separate follow-up (E4 follow-up~L3, now ripe across two
files).

\section{Proof strategy}

The proof is substitute-then-DCT, identical in structure to E4/L1
modulo cast bookkeeping. There are four pieces.

\paragraph{The substitution.} Set $a := t^{1/(2k)}$, so $a^{2k} = t$.
The change of variables $u = a\cdot w$ via Mathlib's
\texttt{Measure.integral\_comp\_mul\_right} reshapes
\[
  t^{1/(2k)} \int_{\mathbb R} e^{-t w^{2k}/(2k)!}\,(\pi(w) - \pi(0))\,dw
\]
into
\[
  \int_{\mathbb R} e^{-u^{2k}/(2k)!}\,(\pi(u/a) - \pi(0))\,du.
\]
The Jacobian bookkeeping is where most of the proof's lines go.
Concretely: $a^{2k} = t$ is proved by routing through
\texttt{Real.rpow\_natCast}, then \texttt{Real.rpow\_mul}, then
\texttt{Real.rpow\_one}, with the intermediate identity
$\alpha \cdot p_{\mathbb R} = 1$ (where $\alpha := 1/p_{\mathbb R}$
and $p_{\mathbb R} := ((2k)_{\mathbb N})_{\mathbb R}$) factored out
as a named \texttt{have} so the calc chain stays linear.

The load-bearing line is the pointwise identity
$(w \cdot a)^{p} = t \cdot w^{p}$. In E4 and L1, this was closed by
\texttt{ring} after specialising to fixed $p \in \{4, 6\}$. For
generic $p = 2k$, \texttt{ring} cannot unify the exponents and
fails. The fix is \texttt{mul\_pow}: $(w\cdot a)^p = w^p \cdot a^p$
unconditionally, then $a^p = t$ via \texttt{ha\_p}, then
\texttt{ring} to commute the multiplication. This was the single
most important tactical correction GPT named at deliberation; without
it the substitution lemma does not close.

\paragraph{The dominator's integrability.} The DCT envelope is
$2M \cdot e^{-u^{2k}/(2k)!}$, which is $t$-independent (the
substitution moved the $t$-dependence into the prior factor). For
generic $k \ge 1$, integrability is established by Gaussian
comparison via the universal pointwise inequality
\[
  x^2 \;\le\; 1 + x^{2k} \qquad (k \ge 1),
\]
proved by case-split on $|x| \le 1$ (where $x^2 \le 1$ and
$x^{2k} \ge 0$) versus $|x| > 1$ (where $x^{2k} \ge x^2$ since
$2k \ge 2$). This gives
$\tfrac{t}{(2k)!} \cdot x^2 - \tfrac{t}{(2k)!} \le \tfrac{t \cdot x^{2k}}{(2k)!}$,
hence $e^{-tx^{2k}/(2k)!} \le e^{t/(2k)!}\cdot e^{-(t/(2k)!) \cdot x^2}$,
the dominator being a Gaussian integrable by Mathlib's
\texttt{integrable\_exp\_neg\_mul\_sq} at $b = t/(2k)!$.

\paragraph{The centred DCT.} The substituted integrand
$F(t,u) := e^{-u^{2k}/(2k)!}\cdot (\pi(u/t^{1/(2k)}) - \pi(0))$ has a
$t$-independent envelope $2M\cdot e^{-u^{2k}/(2k)!}$ (integrable from
the previous step), and its pointwise limit is zero: as $t \to \infty$,
$u/t^{1/(2k)} \to 0$ (the
\texttt{tendsto\_rpow\_atTop} step at exponent $1/(2k) > 0$), and
continuity of $\pi$ at $0$ gives
$\pi(u/t^{1/(2k)}) - \pi(0) \to 0$. Mathlib's filter version of the
dominated convergence theorem\footnote{%
\texttt{tendsto\_integral\_filter\_of\_dominated\_convergence}.}
closes.

\paragraph{The closed-form partition.} The partition
$Z_k(t) = \int e^{-tw^{2k}/(2k)!}\,dw$ has even integrand, so it
equals $2 \cdot \int_{(0,\infty)} e^{-tw^{2k}/(2k)!}\,dw$ via
\texttt{integral\_comp\_abs}. The half-line integral is given by
Mathlib's master half-line moment formula\footnote{%
\texttt{integral\_rpow\_mul\_exp\_neg\_mul\_rpow}.} at
$q = 0$, $p = 2k$, $b = t/(2k)!$, yielding
$(t/(2k)!)^{-1/(2k)} \cdot (1/(2k)) \cdot \Gamma(1/(2k))$. Doubling
gives $(t/(2k)!)^{-1/(2k)} \cdot (2/(2k)) \cdot \Gamma(1/(2k))$, and
$2/(2k) = 1/k$. The exponent flip
$(t/(2k)!)^{-1/(2k)} = ((2k)!/t)^{1/(2k)}$ uses \texttt{inv\_rpow}
and \texttt{Real.rpow\_neg}.

\paragraph{The asymptotic corollary.} Splits
$\pi(w) = (\pi(w) - \pi(0)) + \pi(0)$ inside the integrand. The
centred piece tends to $0$ (by the centred DCT theorem); the
constant piece evaluates via \texttt{kth\_partition} to
$\pi(0)\cdot \tfrac{1}{k}\cdot ((2k)!/t)^{1/(2k)}\cdot \Gamma(1/(2k))$;
multiplying by $t^{1/(2k)}$ and using
$t^{1/(2k)} \cdot ((2k)!/t)^{1/(2k)} = ((2k)!)^{1/(2k)}$ closes.

\section{Roadblocks and resolutions}

This was a single-session formalisation with one substantive
roadblock cluster, all related to the change in tactical class
between fixed-exponent and parametric-exponent proofs.

\paragraph{Cast thrashing in the substitution lemma.} The L1 proof
of $a^p = t$ at fixed $p = 6$ was a calc chain of about ten lines
mixing \texttt{Real.rpow\_natCast}, \texttt{Real.rpow\_mul}, and a
final \texttt{norm\_num}. Generic-$k$ added a layer:
$\alpha \cdot p_{\mathbb R} = 1$ now requires unfolding the
\texttt{set} abbreviations (\texttt{hα\_def}, \texttt{hpR\_def}) so
that \texttt{field\_simp} sees the underlying expressions. The fix
was to factor $\alpha \cdot p = 1$ out as a named \texttt{have}
\emph{before} the calc chain, with explicit \texttt{rw [hα\_def, hpR\_def]} to bring the
unfolded form into scope.

\gptbox{For \texttt{a\textasciicircum{}(2*k) = t}, avoid rewriting to
\texttt{2 * (k : ℝ)} unless you must. The clean route is: prove
\texttt{a \textasciicircum{} p = t} via
\texttt{Real.rpow\_natCast}, then \texttt{Real.rpow\_mul} with
\texttt{α * pR = 1} factored out as a separate \texttt{have}.}

\paragraph{\texttt{ring} cannot unify under generic exponents.} The
load-bearing line in the substitution lemma is
$(w\cdot a)^p = t \cdot w^p$. In E4/L1, both files closed this by
\texttt{ring} after specialising $p$ to $4$ or $6$. For generic
$p = 2k$, \texttt{ring} cannot unify the exponents — it sees $p$ as
an opaque natural number — and fails with a generic ``unsolved
goal''.

\gptbox{In the substitution lemma, use \texttt{mul\_pow}, not
\texttt{ring}. This matters. Fixed exponents like \texttt{4} and
\texttt{6} were happy with \texttt{ring}, but generic
\texttt{(2*k)} won't be.}

\noindent The fix is one line: \texttt{rw [mul\_pow, ha\_p]; ring}.
\texttt{mul\_pow} reshapes $(w\cdot a)^p$ to $w^p \cdot a^p$
unconditionally; \texttt{ha\_p} replaces $a^p$ with $t$;
\texttt{ring} closes the resulting $w^p \cdot t = t \cdot w^p$.

\paragraph{\texttt{set} versus \texttt{exact\_mod\_cast}.}
After \texttt{set fac := (Nat.factorial p : ℝ)}, proving
$0 < \texttt{fac}$ via a direct
\texttt{exact\_mod\_cast Nat.factorial\_pos \_}
fails: the cast tactic sees \texttt{fac} as opaque
and does not unfold the \texttt{set}. The fix is to first
\texttt{change} the goal to the underlying form
\texttt{(0 : ℝ) < (Nat.factorial p : ℝ)}, then dispatch with the
cast. Three sites needed this fix; all minor but each blocked
progress for ${\sim}2$~min while diagnosed.

\paragraph{Negative-exponent rewriting.} Mathlib's
\texttt{integral\_rpow\_mul\_exp\_neg\_mul\_rpow} at $q=0$ produces
the exponent $-(0+1)/p_{\mathbb R}$, parsed as
$(-(0+1))/p_{\mathbb R} = -(1/p_{\mathbb R})$ by Lean's
arithmetic-precedence rules. The single $\mathtt{hexp\_one}$
rewrite covers the positive form $(0+1)/p_{\mathbb R} = 1/p_{\mathbb R}$
that appears in the $\Gamma$ argument, but does \emph{not} fire on
the negative form because the syntactic structure is different.
Adding a parallel $\mathtt{hexp\_neg}$ for the negated form fixed it.

\paragraph{\texttt{set}-defeq vs.\ \texttt{convert} cycles.} The
asymptotic theorem's final goal massage initially used
\texttt{convert ... using 1; funext t; rw [hp\_eq]} chains to
bridge the abbreviated form (using \texttt{p}, \texttt{α},
\texttt{fac} from \texttt{set}) and the public-form goal (using
\texttt{2*k}, the explicit casts, the explicit factorial). Multiple
``no goals to be solved'' errors surfaced because
\texttt{set} makes the abbreviated forms \emph{definitionally equal}
to the public ones, and \texttt{convert} closed the goal at a level
above where I expected. The fix was to rebuild the centred-piece
\texttt{Tendsto} with the \texttt{set} abbreviations directly in
its statement, then use \texttt{kth\_jfunction\_centered\_tendsto\_zero}
\emph{without conversion} (Lean accepts it via defeq), then the
final \texttt{Tendsto.congr'} step closes the public-form goal also
via defeq. Net: about 25~lines of \texttt{convert}/\texttt{funext}
collapsed to 6~lines.

\section{What was Mathlib, what was new}

\paragraph{Mathlib provided.} The full asymptotic stack:
\begin{itemize}[leftmargin=1em,itemsep=0.1em,topsep=0.1em]
\item the master half-line moment formula
  (\texttt{integral\_rpow\_mul\_exp\_neg\_mul\_rpow});
\item the even-function full-line doubling lemma
  (\texttt{integral\_comp\_abs});
\item the substitution-by-constant integral lemma
  (\texttt{Measure.integral\_comp\_mul\_right});
\item the filter version of dominated convergence%
  \footnote{\texttt{tendsto\_integral\_filter\_of\_dominated\_convergence}.};
\item the rpow-tends-to-infinity step (\texttt{tendsto\_rpow\_atTop});
\item integrability of the Gaussian dominator
  (\texttt{integrable\_exp\_neg\_mul\_sq}).
\end{itemize}
The $\Gamma$-function side is \texttt{Real.Gamma} from
\texttt{Mathlib.Analysis.SpecialFunctions.Gamma.Basic}.

\paragraph{What was new.} Three things, in increasing order of
likely future reuse:

\begin{itemize}[leftmargin=1em,itemsep=0.2em]
\item \texttt{sq\_le\_one\_add\_pow\_two\_mul}: the universal pointwise
  inequality $x^2 \le 1 + x^{2k}$ for $k \ge 1$. Cleaner and more
  general than the L1 file's bespoke
  \texttt{sq\_le\_one\_add\_pow\_six} or the Quartic file's
  square-completion identity. Likely candidate for promotion to
  \texttt{lean-common} once a second use materialises.
\item \texttt{kth\_integrable}: full-line integrability of
  $e^{-t x^{2k}/(2k)!}$ for $k \ge 1$, via the universal Gaussian
  comparison above. Generalises and unifies the existing
  fixed-exponent integrability helpers in the quartic and sextic
  files at $n = 0$.
\item \texttt{kth\_partition}: closed-form full-line monomial
  partition function. Generalises the existing fixed-exponent
  partitions for the quartic and sextic potentials (which remain
  valid as standalone theorems on the existing branches; once E4 and
  L1 merge to \texttt{main}, they become 5-line specialisations,
  and the cleanup pass is a follow-up tide).
\end{itemize}

\noindent The substitution lemma \texttt{kth\_jfunction\_centered\_subst}
and the centred DCT theorem are private and bespoke; they would not
add value to \texttt{lean-common} without further generalisation
(e.g.\ to \emph{any} potential admitting a power-law substitution).

\section{Lessons}

\begin{itemize}[leftmargin=1em,itemsep=0.2em]
\item \emph{For \texttt{laplace/CLAUDE.md}:} \texttt{ring} cannot
  unify exponents involving free natural-number variables. For
  $(w\cdot a)^p$ with parametric $p$, use \texttt{mul\_pow}
  explicitly, not \texttt{ring}. The quartic and sextic files use
  \texttt{ring} as a fixed-exponent special case; generic-$p$
  proofs must replace it.

\item \emph{For \texttt{laplace/CLAUDE.md}:} \texttt{exact\_mod\_cast}
  does \emph{not} unfold \texttt{set} abbreviations. After
  \texttt{set fac := (Nat.factorial p : ℝ)}, prove $0 < \texttt{fac}$
  with an explicit \texttt{change} to the underlying form before
  the cast. (Already in \texttt{CLAUDE.md} as a related gotcha for
  cast directions; this run extends it to the \texttt{set}-defeq
  case.)

\item \emph{For \texttt{laplace/CLAUDE.md}:} when a Mathlib lemma
  produces an exponent of the form $-(q+1)/p$ in the conclusion (as
  \texttt{integral\_rpow\_mul\_exp\_neg\_mul\_rpow} does), Lean
  parses it as $(-(q+1))/p$ — distinct from $-((q+1)/p)$ for
  rewriting purposes. Provide \emph{both} the positive-form rewrite
  (for the $\Gamma$ argument) and the negative-form rewrite (for the
  rpow exponent) when massaging.

\item \emph{For the tide skill:} the deliberation transcript carried
  forward all three load-bearing tactical points (\texttt{mul\_pow},
  the \texttt{set}/\texttt{change} pattern, the explicit-cast
  preference) into the proof without any mid-proof
  thrashing-rescue consult needed. This is the second tide in a row
  (Tide~I4 was the first) where the GPT-5.5~Pro consult at Step~2
  contained \emph{enough} tactical detail that Step~3 was mechanical
  Lean idiom. The pattern: when Step~2 has already named the
  Mathlib idiom and the specific friction points, the proof writes
  itself.

\item \emph{For \texttt{tide-pick}:} the L1 retrospective named N1
  as ``now ripe'' \emph{within hours} of L1 landing. Tide-pick
  picked it as the top recommendation in the v3 survey within
  another hour. The pattern of ``retrospective Follow-ups become
  next-tide candidates'' continues to be the dominant source of
  high-value tide steps; the strict-improvement-of-the-day pipeline
  is producing single-session formalisations at sub-day cadence.
\end{itemize}

\section{Follow-ups}

\begin{itemize}[leftmargin=1em,itemsep=0.2em]
\item \emph{Refactor \texttt{JFunctionQuartic.lean} and
  \texttt{JFunctionSextic.lean} to be 5-line specialisations of
  \texttt{kth\_jfunction\_asymptotic}.} Cannot be done in this Tide
  because both files are on unmerged branches at base time. After
  E4, L1, and N1 all land on \texttt{main}, this is a focused
  cleanup pass: each existing file shrinks from ${\sim}280$~lines
  to ${\sim}30$~lines (the specialisation plus constant unification
  for the user-facing theorem name). One follow-up tide for the pair.

\item \emph{Hypothesis weakening across all three files.} The
  deliberation's mathematically-minimal hypothesis class is
  \texttt{AEMeasurable + ContinuousAt 0 + global bound}; all three
  files use the slightly stronger \texttt{Continuous prior + global bound}.
  Now that the parametric form exists, weakening should land in N1
  (\texttt{kth\_*}) only — the specialisations inherit. Expected
  ${\sim}50$~lines of \texttt{comp\_quasiMeasurePreserving} plumbing.

\item \emph{Promote \texttt{sq\_le\_one\_add\_pow\_two\_mul} to
  \texttt{lean-common}.} Genuinely general; subsumes the L1 file's
  bespoke \texttt{sq\_le\_one\_add\_pow\_six}. Defer until a second
  downstream use (probably the hypothesis-weakening tide above).

\item \emph{Promote \texttt{kth\_partition} to \texttt{lean-common}.}
  The natural full-line monomial partition formula; clean addition.
  Defer until a second use; the most likely trigger is the
  bounded-prior-monomial generalisation tide (B2 in the survey),
  which would consume it as primitive infrastructure.

\item \emph{Asymptotic-equivalence wrapper $J_k(t) \sim \pi(0)\cdot Z_k(t)$
  (E4 follow-up~L2, generalised).} A 15-line corollary that takes
  the \texttt{Tendsto (t\textsuperscript{1/(2k)} · J\textsubscript{k})}
  form and divides through by \texttt{kth\_partition}. Now applies
  to all $k \ge 1$ in one shot rather than once per file. Smallest
  candidate on the board going forward.

\item \emph{\texttt{\textbackslash leanref}-tag the grammar paper.}
  Once the grammar paper has TeX statements pinning all three
  $k = 2, 3, $ generic-$k$ J-function asymptotics, tag with
  \texttt{\textbackslash leanref} pinned to commit
  \texttt{ddb1f14} (this tide). Process tide; depends on
  paper-side draft.
\end{itemize}

\section*{Acknowledgements}

GPT-5.5~Pro contributed the deliberation-stage advisory at Step~2,
including the three tactical points that made the proof
single-session: the \texttt{mul\_pow} recipe, the
$\alpha \cdot p = 1$ factoring pattern, and the explicit-cast
preference for the public theorem statement. The full consult is
preserved at
\texttt{projects/primer/tide-log/gpt\_responses\_n1/n1\_v1.md}. No
mid-proof thrashing-rescue consult was needed; the Step~2 advisory
was sufficient.

\end{document}
